
\section*{Discussion}

\begin{figure}
	\centering
	\fbox{\rule[-.5cm]{4cm}{4cm} \rule[-.5cm]{4cm}{0cm}}
	\caption{Sample figure caption.}
	\label{fig:fig1}
\end{figure}

\begin{table}
	\caption{Sample table title}
	\centering
	\begin{tabular}{lll}
		\toprule
		\multicolumn{2}{c}{Part}                   \\
		\cmidrule(r){1-2}
		Name     & Description     & Size ($\mu$m) \\
		\midrule
		Dendrite & Input terminal  & $\sim$100     \\
		Axon     & Output terminal & $\sim$10      \\
		Soma     & Cell body       & up to $10^6$  \\
		\bottomrule
	\end{tabular}
	\label{tab:table}
\end{table}

\begin{itemize}
	\item Lorem ipsum dolor sit amet
	\item consectetur adipiscing elit.
	\item Aliquam dignissim blandit est, in dictum tortor gravida eget. In ac rutrum magna.
\end{itemize}


** In silico is not the same as biological validation; Lack of wet lab evidence is a weakness of this paper

Does RNA seq data present evidence for expression? Weakness of that argument.
But, single time point, single stresser.

Uncertain if sp_algae is expressed or where it may go to.

** Unknown metal specificity

Metal specificity is a spectrum, not a binary: there are Mn-specific,
Fe-specific, and cambialistic SODs that work with either ion, and the modern
distribution is polyphyletic. Specificity is tuned by the redox potential of
the active-site metal, and several non-conserved second-sphere residues
contribute to that tuning; consequently prediction methods based solely on
primary sequence have not given strong predictive values. There are documented
sequence/metal mismatches — e.g. the Sinorhizobium meliloti SOD has an
iron-type sequence but preferentially incorporates Mn in vivo. The mechanistic
anchor most methods lean on is the position of the solvent-coordinating
glutamine (α2-helix in typical FeSODs vs the C-terminal end of β2 in most
MnSODs), but site-directed swaps show this is not the only determinant of metal
selection. There is a model - SODa - but not super predictive. The lineage of
this work: Parker & Blake (1988) first showed Fe- and Mn-SODs can be
distinguished by primary structure; Wintjens et al. (2004, JBC) and Wintjens,
Gilis & Rooman (2008, Proteins) formalized the fingerprints. A simpler,
hand-usable heuristic also exists — cyanobacterial SODs can be split by a
three-trait key, e.g. a signature where FeSOD's F184/A188/Q189…T280…F/Y303 is
replaced in MnSOD by R184/G188/G189…G280…W303 (positions are
cyanobacteria-specific).

The metal is set by the active-site environment of the protein, and
localization is set by the presequence — two independent properties.
Localization cannot pin the functional metal at all in the ambiguous cases.

Mn supplementation may be good https://www.nature.com/articles/s41598-018-34994-4

Mn dependence https://journals.plos.org/climate/article?id=10.1371/journal.pclm.0000897

Discusses Mt/S and Chloroplast versions in algae, but they took annotations
from published sources and presented alignments assuming those annotations were
correct. They noticed bigger differences between their MnSOD3 group with other
groups. This is because their MnSOD3 is in fact a SP targeting one, while their
MnSOD1 and MnSOD2 groups are actually our sp_cpt groups.

Encoded protein sequence length within each MnSOD isoform was similar across
clades (within 1–9 aa of each other), but absolute length between isoforms
varied by up to 19% (e.g., SymMnSOD3 vs. SymMnSOD2). Mean pairwise protein
identities across clades were highest for SymFeSOD (96% ± 2%; mean ± SD),
SymMnSOD1 (94% ± 2%), and SymMnSOD2 (90% ± 3%), with near identical amino acid
sequences (96-97%) for MnSOD1 within clade C (C1, C3, C15) (Figure 1,
Additional file 1). The predicted tertiary monomer structure of SymMnSOD1 from
Symbiodinium B1 and C1 was also highly similar (Figure 2).

C15_M_digitata_SymMnSOD1_KJ6 -            265 sp_ctp_MnFeSOD_algae_mature -            193  1.8e-140  456.0   2.3   1   1  2.2e-140  2.2e-140  455.7   2.3     1   193    58   250    58   250 1.00 -
C3_Mp_SymMnSOD1_KJ672520     -            265 sp_ctp_MnFeSOD_algae_mature -            193  1.9e-140  455.9   2.0   1   1  2.3e-140  2.3e-140  455.7   2.0     1   193    58   250    58   250 1.00 -
D_A_hyacinthus_SymMnSOD2_GAF -            294 sp_ctp_MnFeSOD_algae_mature -            193  4.1e-140  454.8   2.1   1   1  5.1e-140  5.1e-140  454.5   2.1     1   193    55   247    55   247 1.00 -
F1_CCMP2468_SymMnSOD2_Assemb -            297 sp_ctp_MnFeSOD_algae_mature -            193  4.6e-140  454.6   2.1   1   1  5.8e-140  5.8e-140  454.3   2.1     1   193    58   250    58   250 1.00 -
B1_Mf1_05b_SymMnSOD2_Assembl -            299 sp_ctp_MnFeSOD_algae_mature -            193  6.2e-140  454.2   1.9   1   1  7.8e-140  7.8e-140  453.9   1.9     1   193    58   250    58   250 1.00 -
B1_Mf1_05b_SymMnSOD2_rep_c13 -            303 sp_ctp_MnFeSOD_algae_mature -            193  6.5e-140  454.2   1.9   1   1  8.1e-140  8.1e-140  453.9   1.9     1   193    58   250    58   250 1.00 -
C1_CCMP2466_SymMnSOD1_AHI543 -            263 sp_ctp_MnFeSOD_algae_mature -            193  1.3e-138  449.9   1.9   1   1  1.5e-138  1.5e-138  449.7   1.9     1   193    56   248    56   248 1.00 -
B1_Ap1_SymMnSOD1_KJ672521    -            261 sp_ctp_MnFeSOD_algae_mature -            193  5.7e-138  447.8   2.4   1   1  6.7e-138  6.7e-138  447.6   2.4     1   193    58   250    58   250 1.00 -
A1_Casskb8_MnSOD__c18645     -            281 sp_ctp_MnFeSOD_algae_mature -            193    6e-138  447.8   2.1   1   1  7.3e-138  7.3e-138  447.5   2.1     1   193    61   253    61   253 1.00 -
B1_Mf1_05b_SymMnSOD1_Assembl -            265 sp_ctp_MnFeSOD_algae_mature -            193  1.2e-137  446.8   2.2   1   1  1.4e-137  1.4e-137  446.5   2.2     1   193    58   250    58   250 1.00 -
A1_Casskb8_MnSOD_rep_c710    -            270 sp_ctp_MnFeSOD_algae_mature -            193  1.1e-136  443.7   1.3   1   1  1.3e-136  1.3e-136  443.4   1.3     1   193    63   255    63   255 1.00 -
A1_CCMP2467_MnSOD_Assembly1  -            193 sp_ctp_MnFeSOD_algae_mature -            193   9.3e-91  293.8   0.8   1   1   1.1e-90   1.1e-90  293.5   0.8     1   133    61   193    61   193 1.00 -
A1_CCMP2467_MnSOD_Assembly4  -            189 sp_ctp_MnFeSOD_algae_mature -            193   1.2e-86  280.3   0.8   1   1   1.5e-86   1.5e-86  280.1   0.8     1   127    63   189    63   189 1.00 -
A1_Casskb8_SymMnSOD3_rep_c41 -            246 sp_ctp_MnFeSOD_algae_mature -            193   4.3e-74  239.4   0.6   1   1   5.1e-74   5.1e-74  239.1   0.6     2   193    24   219    23   219 0.93 -
D_A_hyacinthus_SymMnSOD3_GAF -            247 sp_ctp_MnFeSOD_algae_mature -            193   7.6e-70  225.5   0.2   1   1     9e-70     9e-70  225.3   0.2     2   193    25   220    24   220 0.94 -
C3_A_aspera_SymMnSOD1_FE8660 -            105 sp_ctp_MnFeSOD_algae_mature -            193     6e-65  209.6   2.6   1   1   6.8e-65   6.8e-65  209.4   2.6   104   193     1    90     1    90 1.00 -
A_PF_2005_SymFeSOD_AY916504  -            201 sp_ctp_MnFeSOD_algae_mature -            193     2e-53  171.9   1.3   1   1   2.4e-53   2.4e-53  171.7   1.3     1   191     4   191     4   193 0.91 -
A1_CCMP2467_MnSOD_Assembly3  -             80 sp_ctp_MnFeSOD_algae_mature -            193   5.9e-45  144.3   0.9   1   1   6.7e-45   6.7e-45  144.1   0.9   129   193     1    65     1    65 0.99 -
E_CCMP421_SymFeSOD_KJ672517  -            149 sp_ctp_MnFeSOD_algae_mature -            193   3.9e-42  135.1   1.9   1   1   4.4e-42   4.4e-42  134.9   1.9    22   174     1   149     1   149 0.90 -
A1_CCMP2467_SymMnSOD3_Assemb -            173 sp_ctp_MnFeSOD_algae_mature -            193   1.7e-39  126.5  14.1   1   2   2.3e-38   2.3e-38  122.8  10.7     2   148    24   156    23   160 0.93 -
A1_CCMP2467_SymMnSOD3_Assemb -            173 sp_ctp_MnFeSOD_algae_mature -            193   1.7e-39  126.5  14.1   2   2     0.083     0.083    3.5   1.1   181   193   152   164   151   164 0.88 -
B1_Ap1_SymFeSOD_KJ672519     -            148 sp_ctp_MnFeSOD_algae_mature -            193   3.7e-38  122.1   0.9   1   1   4.4e-38   4.4e-38  121.9   0.9     7   149     1   143     1   147 0.90 -
B1_Mf1_05b_MnSOD__c29099     -            148 sp_ctp_MnFeSOD_algae_mature -            193   2.2e-37  119.6   0.2   1   1   2.6e-37   2.6e-37  119.4   0.2     2   124    20   147    19   148 0.93 -
F1_Mv_SymFeSOD_KJ672518      -            132 sp_ctp_MnFeSOD_algae_mature -            193   5.7e-35  111.7   1.5   1   1   6.2e-35   6.2e-35  111.6   1.5    20   149     2   131     1   132 0.90 -
A1_CCMP2467_SymMnSOD3_Assemb -             86 sp_ctp_MnFeSOD_algae_mature -            193   2.2e-27   87.0   0.3   1   1   2.6e-27   2.6e-27   86.7   0.3   137   193     1    59     1    59 0.94 -
B1_Mf1_05b_SymMnSOD3__c19510 -             86 sp_ctp_MnFeSOD_algae_mature -            193   2.2e-27   87.0   0.3   1   1   2.6e-27   2.6e-27   86.7   0.3   137   193     1    59     1    59 0.94 -
A1_Casskb8_SymMnSOD3_rep_c41 -            246 sp_MnFeSOD_algae_mature -            197  5.8e-133  431.3   2.4   1   1  6.7e-133  6.7e-133  431.1   2.4     1   197    23   219    23   219 1.00 -
D_A_hyacinthus_SymMnSOD3_GAF -            247 sp_MnFeSOD_algae_mature -            197  2.7e-131  425.9   0.6   1   1  3.2e-131  3.2e-131  425.7   0.6     1   197    24   220    24   220 1.00 -


** Localization

We checked how DeepLoc and SignalIP would call alveolata proteins that
Should be mitochondria, using K00235 as control: DeepLoc detected mTP.  Should
be secreted, using K09490 (ER chaperone) as control: sequences from various
species fall into two groups, one had SP+ER DeepLoc call, the other
NLS+Cytoplasm. Stripping out the SP, DeepLoc does not by default report Plastid
but reports ER and Cytoplasm as localization, and NLS, NES, and SP as signals.
Should be plastid, using K02639 (ferredoxin) as control: DeepLoc detected Golgi
as localization, but SP was only called for some sequences not all even though
the leader sequence appears to be highly conserved. Stripping out the SP,
DeepLoc reports Plastid as the localization for the truncated protein. n is
very small here, as many of the sequences are essentially the same 3 protein
sequences. To make this publishable we should look for different plastid
targeted sequences and hope to observe the same behavior from DeepLoc.

** HMM profiles pan metagenomic or family specific

Because K04564 encapsulates the entire Fe-Mn SOD family across all domains of
life (Bacteria, Archaea, and Eukaryota) [Source: KEGG Orthology K04564], a high
bit-score threshold (e.g., >150) would optimize precision for eukaryotic
sequences but drastically collapse the recall rate for highly divergent
bacterial or archaeal orthologs. The optimization algorithm drove the threshold
down to 41.53 to maintain functional annotation capacity across highly diverse
phyla.

In an ancient, universally distributed enzyme like MnSOD, large stretches of
the ~200 amino acid sequence (such as surface-exposed loop regions and variable
structural linkers) exhibit extreme sequence divergence across species. In the
pHMM, these hypervariable positions have high entropy, dropping their bit-score
contribution to near zero. Only the rigid structural core—specifically the
N-terminal alpha-hairpin domain, the C-terminal alpha/beta domain, and the 4
metal-coordinating catalytic nodes (His-50, His-98, Asp-183,
His-187)—contribute highly positive bits. A truly orthologous but highly
divergent sequence will only accumulate score points at these narrow, invariant
motifs, resulting in a low overall sequence bit score that still represents a
true positive.

E.g. even in Alveolata, can have quite different sequences, so profiles may
need different thresholds. Hard to quantify though if they are differentiating.

Using HMM profiles with rules to check for leaders and key residues. Profiles
give you a fast first screen, and then the key residues make sure they are
catalytically functional. Leaders verify full length proteins. Mature HMM gives
you a quick way to see what class a sequence may be w/o detecting a leader
which could be challenging esp in algae.

** Considerations for RNAseq analysis

In addition to more specific HMM profiles to match to, there is the issue of
sequencing holobiont and then doing capture or removal

** Continuing work of curating proteins

n is small, as we get more data, the profiles get better and new profiles
likely emerge. We are also interested in taking the approach to other proteins
and protein families.

Gene model is difficult, needs more RNAseq data, AUGUSTUS modifications do not
work. HMM of the mature protein is useful for global HMM search.


